\section{The four-term identity}\label{sec:local}

This section contains the only local identity used in the proof, together with
two facts about highest reachable vertices.
We state them for directed graphs rather than orientations, since later we
apply them to graphs obtained by adding or deleting single edges.

Let $D$ be a directed graph on $[n]$ whose edges point upwards, or a periodic
directed graph on $\setZ$ whose edges point upwards and whose projection to
$\setZ/n\setZ$ is acyclic.
In both cases every vertex $z$ has a largest reachable vertex, which we denote by
$H_D(z)$; in the periodic case, the bound follows as in
Section~\ref{subs:circularDefinitions}.
The function $H_D$ is weakly decreasing along every edge of $D$, and $H_D(z)$ is
the largest sink of $D$ reachable from $z$.

\begin{lemma}\label{lem:ineffective}
Let $D$ be as above.
\begin{enumerate}[(i)]
\item Let $F$ be a set of upward pairs, translation invariant in the periodic
case, such that $H_D(y)\leq H_D(x)$ for every $xy\in F$.
Then $H_{D\cup F}=H_D$, and in the periodic case $D\cup F$ still has acyclic
projection.
\item If $D$ contains a path from $x$ to $y$ avoiding the edge $xy$, then
deleting $xy$ does not change reachability in $D$.
\end{enumerate}
\end{lemma}

\begin{proof}
For (i), $H_D$ is weakly decreasing along every edge of $D\cup F$.
Hence every vertex reachable from $z$ in $D\cup F$ is at most $H_D(z)$, and the
maximizing path in $D$ is still available.
In the periodic case, a directed cycle in the projection would lift to a path
from $z$ to $z+kn$ for some $k\geq1$, and then
$H_D(z)+kn=H_D(z+kn)\leq H_D(z)$, which is absurd.
Part (ii) is clear.
\end{proof}

\begin{lemma}\label{lem:sinkRelabel}
Let $D$ and $D'$ be directed graphs on $[n]$ with upward edges, and let
$\sigma$ be a permutation of $[n]$ mapping the edges of $D$ bijectively onto the
edges of $D'$.
If $\sigma$ is increasing on the set of sinks of $D$, then
$H_{D'}(\sigma(z))=\sigma(H_D(z))$ for all $z$.
In particular the fibres of $H_{D'}$ are the images of the fibres of $H_D$.
\end{lemma}

\begin{proof}
The vertex $H_D(z)$ is the largest sink of $D$ reachable from $z$, and $\sigma$
maps the sinks reachable from $z$ in $D$ onto the sinks reachable from
$\sigma(z)$ in $D'$, preserving their order.
\end{proof}

We say that an edge is \defin{ineffective} in $D$ if it satisfies the hypothesis
of Lemma~\ref{lem:ineffective}(i).

\begin{lemma}[The four-term identity]\label{lem:fourTerm}
Let $D$ be a directed graph on $[n]$ with upward edges.
Let $u<v<w$ be vertices such that $uw$ is an edge of $D$, and $uv$ and $vw$ are
not.
Put $D_{00}=D$, $D_{10}=D+uv$, $D_{01}=D+vw$ and $D_{11}=D+uv+vw$, and write
$\Kp(D_{ij})$ for $\Kp$ of the transitive closure.
\begin{enumerate}[(i)]
\item For every coloring $\kappa:[n]\to\setNN$ which increases along the edges
of $D$,
\[
 [\kappa \text{ increases on } D_{10}]+[\kappa \text{ increases on } D_{01}]
 =1+[\kappa \text{ increases on } D_{11}].
\]
In particular $\Kp(D_{10})+\Kp(D_{01})=\Kp(D_{00})+\Kp(D_{11})$, and the same
holds for the sums over colorings with prescribed colors at any set of vertices.
\item The four graphs can be paired, either as $\{D_{00},D_{01}\}$ and
$\{D_{10},D_{11}\}$ or as $\{D_{00},D_{10}\}$ and $\{D_{01},D_{11}\}$, so that
paired graphs have the same function $H$.
In particular $\Phi(D_{10})+\Phi(D_{01})=\Phi(D_{00})+\Phi(D_{11})$ for every
function $\Phi$ of the highest reachable vertex map, for example
$\Phi=\elementaryE_{\lambda}$ of its fibre sizes.
\end{enumerate}
\end{lemma}

The lemma has no cover hypotheses, and the other edges of $D$ are arbitrary.
In the diagram notation of~\cite{AlexanderssonSulzgruber2022}, the identity
reads
\[
\ytableausetup{boxsize=1.1em}
\begin{ytableau}
\to & 0 & *(yellow) w \\
1 & *(yellow) v \\
*(yellow) u
\end{ytableau}
\;+\;
\begin{ytableau}
\to & 1 & *(yellow) w \\
0 & *(yellow) v \\
*(yellow) u
\end{ytableau}
\;=\;
\begin{ytableau}
\to & 0 & *(yellow) w \\
0 & *(yellow) v \\
*(yellow) u
\end{ytableau}
\;+\;
\begin{ytableau}
\to & 1 & *(yellow) w \\
1 & *(yellow) v \\
*(yellow) u
\end{ytableau}\,,
\]
where $1$ and $0$ indicate whether a short edge is present, the long edge is
present in all four graphs, and all other edges are the same throughout.

\begin{proof}
For (i), write $a,b,c$ for the colors at $u,v,w$, so that $a<c$.
Then $[a<b]+[b<c]=1+[a<b<c]$, which is the displayed identity.
Summing $x_\kappa$ over all $\kappa$ increasing on $D$, with any colors
prescribed, gives the second statement.

For (ii), let $H=H_D$.
Since $uw$ is an edge, $H(w)\leq H(u)$.
Suppose first that $H(w)\leq H(v)$.
Then $vw$ is ineffective in $D$.
Adding $uv$ does not change $H$ at $v$ or at $w$, since no path starting at
$v$ or $w$ uses an edge leaving $u<v$.
Hence $vw$ is also ineffective in $D_{10}$, and Lemma~\ref{lem:ineffective}
pairs $D_{00}$ with $D_{01}$ and $D_{10}$ with $D_{11}$.
Suppose instead that $H(v)<H(w)$.
Then $uv$ is ineffective in $D$, since $H(v)<H(w)\leq H(u)$.
In $D_{01}$, a path from $v$ either avoids $vw$ or continues from $w$, so
$H_{D_{01}}(v)\leq \max(H(v),H(w))=H(w)\leq H_{D_{01}}(u)$.
Thus $uv$ is ineffective in $D_{01}$ as well, and we obtain the second pairing.
\end{proof}

Lemma~\ref{lem:fourTerm} does not say that the two middle terms correspond to
each other.
In the first case of the proof, $D_{10}$ has the same highest reachable vertices
as $D_{11}$, not as $D_{01}$.
This is why the identity is used for sums over many orientations at once.
